\documentclass[lettersize,journal]{IEEEtran}
\usepackage{amsmath,amsfonts}
%\usepackage{algorithmic}
\usepackage{algorithm}
\usepackage{algpseudocode}
\usepackage{array}
\usepackage[caption=false,font=normalsize,labelfont=sf,textfont=sf]{subfig}
\usepackage{textcomp}
\usepackage{stfloats}
\usepackage{url}
\usepackage{verbatim}
\usepackage{graphicx}
\usepackage{cite}
\usepackage{booktabs} % For more aesthetically pleasing table lines
\usepackage{multirow} % Required for merging rows
\usepackage{pgfplotstable}
\usepackage{makecell}

\algrenewcommand\algorithmicrequire{\textbf{Input:}}
\algrenewcommand\algorithmicensure{\textbf{Output:}}
\algrenewcommand\algorithmiccomment[1]{\hfill/* #1 */}

\usepackage{tabularx} % 导言区加入
\usepackage[table]{xcolor}
\usepackage{tcolorbox}
\newcommand{\red}[1]{\textcolor{red}{#1}}
\newcommand{\best}[1]{\cellcolor{gray!65}\textbf{#1}}
\newcommand{\second}{\cellcolor{gray!45}}
\newcommand{\third}{\cellcolor{gray!20}}

\newif\ifshowcmt
\showcmttrue   % 显示批注
% \showcmtfalse % 隐藏批注

\newcommand{\cmt}[1]{%
	\ifshowcmt
	\textcolor{blue}{/* #1 */}%
	\fi
}

\hyphenation{op-tical net-works semi-conduc-tor IEEE-Xplore}
% updated with editorial comments 8/9/2021

\begin{document}

\title{MetTrain: A Metamorphic-Relation-Driven Training Framework for Natural Language Inference}

\author{Zhehao Li, Jinfu Chen, Tsong Yueh Chen, Saihua Cai, Chengying Mao, Zhicheng Wu
        % <-this % stops a space
%\thanks{This paper was produced by the IEEE Publication Technology Group. They are in Piscataway, NJ.}% <-this % stops a space
\thanks{Zhehao Li, Jinfu Chen, Saihua Cai, Zhicheng Wu are with School of Computer Science \& Communication Engineering, Jiangsu University, Zhenjiang, 212013, China and also with the Jiangsu Key Laboratory of Security Technology for Industrial Cyberspace, Zhenjiang 212013, China. (e-mail: zhehaoli@stmail.ujs.edu.cn, jinfuchen@ujs.edu.cn, caisaih@ujs.edu.cn, zhichengwu@stmail.ujs.edu.cn)(\textit{Corresponding author: Jinfu Chen.}).}
\thanks{Tsong Yueh Chen is with \red{Department of Computing Technologies}, Swinburne University of Technology, Hawthorn, VIC 3122, Australia (e-mail: tychen@swin.edu.au).}
\thanks{Chengying Mao is with the School of Software and IoT Engineering, Jiangxi University of Finance and Economics, Nanchang 330013, China (e-mail: maochy@yeah.net).}}

% The paper headers
%\markboth{Journal of \LaTeX\ Class Files,~Vol.~14, No.~8, April~2026}%
%{Shell \MakeLowercase{\textit{et al.}}: A Sample Article Using IEEEtran.cls for IEEE Journals}

%\IEEEpubid{0000--0000/00\$00.00~\copyright~2021 IEEE}
% Remember, if you use this you must call \IEEEpubidadjcol in the second
% column for its text to clear the IEEEpubid mark.

\maketitle

\begin{abstract}
Improving the reliability of natural language inference (NLI) models remains an open challenge, especially under limited supervision and distribution shifts. \cmt{This is a standard term in AI, it refers to a mismatch between the training and test data distributions. For example, training an NLI model on one dataset, such as SICK, and evaluating it on another dataset, such as SNLI, constitutes a cross-dataset distribution shift or out-of-domain setting.} Existing semi-supervised learning methods mainly rely on confidence-based pseudo-labeling generated by a fine-tuned LLM, which often fails to enforce logical consistency or maintain robustness under distribution shifts. To alleviate this, we propose MetTrain, a metamorphic-relation-driven semi-supervised learning framework for NLI. Unlike conventional uses of metamorphic relations (MRs) as testing criteria for checking system behavior under controlled input transformations, MetTrain uses them to guide training-time data generation. Using 14 MRs, including composite MRs, MetTrain generates premise--hypothesis--label triples from labeled and unlabeled text via a lightweight 4B-parameter language model, and trains the target NLI model with R-Drop regularization. Experiments on multiple datasets show that MetTrain \red{achieves consistent but moderate in-domain gains} over 17 comparison methods, \red{while showing larger gains under out-of-domain (OOD) evaluation}, with \red{absolute} improvements of up to 1.93 percentage points in average accuracy and 7.00 \red{percentage points in macro-F1 under cross-dataset OOD transfer, e.g., when training on SICK and evaluating on SNLI.} The results suggest that MRs can both expand supervision through metamorphic synthesis and improve robustness by enforcing explicit logical constraints.
\end{abstract}

\begin{IEEEkeywords}
Metamorphic testing, Metamorphic relation, Semi-supervised learning, Natural language inference
\end{IEEEkeywords}

\section{Introduction}

``No man's knowledge here can go beyond his experience.'' — John Locke \cite{locke1847essay}. 

Deep neural networks inherit this empiricist constraint: they achieve state-of-the-art (SOTA) performance by ingesting massive, meticulously annotated corpora \cite{DBLP:conf/emnlp/BowmanAPM15}, which serve as the digital analogue of experience. This dependency is particularly acute in Natural Language Inference (NLI), which requires large logic-annotated datasets to determine whether one sentence entails another. Yet, collecting such datasets at scale remains time-consuming and expensive \cite{nie-etal-2020-adversarial, DBLP:conf/acl/VarshneyBGB22,DBLP:conf/naacl/WilliamsNB18}, motivating techniques that reduce reliance on labeled data.

Semi-supervised learning (SSL) alleviates this problem by leveraging unlabeled data, but it faces a critical bottleneck: ensuring the logical validity and semantic diversity of generated supervision signals. While existing approaches typically rely on data augmentation (DA)~\cite{DBLP:journals/asc/PellicerFC23} or pseudo-labeling (PL)~\cite{lee2013pseudo}, they face significant trade-offs. Rule-based techniques are resource-efficient but often yield only incremental gains due to limited semantic diversity, while model-based generation can introduce richer variation at higher cost~\cite{feng-etal-2021-survey}. Recent SOTA methods employ teacher--student-like pipelines~\cite{DBLP:conf/emnlp/SadatC22a,DBLP:conf/emnlp/ParkC24}, but they typically require costly fine-tuning. Such methods are often bounded by the teacher model's capability, and may offer limited out-of-domain generalization~\cite{DBLP:journals/corr/abs-2202-07136}. When the test distribution drifts, even high-performing models can fail~\cite{mccoy-etal-2019-right, gururangan-etal-2018-annotation}, revealing their reliance on shallow heuristics rather than the inference logic that robust NLI requires.

To bridge this gap, we revisit SSL for NLI through the lens of \emph{Metamorphic Testing} (MT)~\cite{DBLP:journals/corr/abs-2002-12543} and propose \textbf{MetTrain}, a metamorphic-relation-driven SSL framework. In traditional software engineering, metamorphic testing uses \emph{Metamorphic Relations} (MRs) as oracles to detect bugs by checking whether outputs satisfy domain-specific properties. More recently, MT has been increasingly adopted to test NLP systems, particularly for assessing robustness under controlled semantic or syntactic transformations~\cite{11185922}. 

\red{Recent reflections on MT further suggest that MRs can play broader roles beyond alleviating the test oracle problem and may bring new concepts to other disciplines~\cite{DBLP:conf/sigsoft/ChenT21}. Inspired by this broader view, we invert the conventional MT paradigm: }instead of using MRs solely for post-hoc validation, we elevate them to \emph{the source of training-time supervision signals}. \red{Specifically, this paper focus on a common special class of MRs, where each MR maps one source NLI instance to one follow-up NLI instance (as defined in Definition 1 of Section III-A).} By treating MRs as explicit ``knowledge carriers''~\cite{DBLP:journals/corr/abs-2104-04718}, MetTrain imposes checkable logical constraints on sample generation. Concretely, we define 14 MRs (including composite MRs) to enforce logical consistency. We separate \emph{inference logic} from \emph{surface realization}: a compact, frozen large language model (LLM) (4B parameters) serves merely as an MR executor to linguistically instantiate the rigid logic defined by the MRs; the synthesized samples are then used to train the target NLI model with R-Drop regularization.

This framework integrates metamorphic synthesis with robustness-oriented training, providing richer supervision while improving robustness and logical consistency under limited labeled data. Importantly, MetTrain achieves these gains with a compact frozen 4B LLM \cmt{Here, ``compact'' means that the generator LLM has only 4B parameters, which is much smaller than commonly used 7B/8B or larger LLMs. ``Frozen'' means that the generator LLM is used without fine-tuning; its parameters are not updated during MetTrain.}, suggesting that the framework does not rely on a large or heavily optimized generator. Empirically, MetTrain shows consistent improvements in both performance and robustness over 17 competitive baselines. These results indicate that metamorphic-relation-driven training offers a practical and promising \red{approach} for improving reliability in NLI, and may also be useful for other NLP tasks that require logically consistent behavior.
\textbf{Our contributions are summarized as follows:}
\begin{itemize}
	\item \textbf{MR-Driven Supervision for Semi-Supervised NLI:} We propose MetTrain, an MR-driven training framework for semi-supervised NLI that extends the role of metamorphic relations from post-hoc test oracles to the source of training-time supervision signals. 
	\item \textbf{Empirical Evaluation:} We evaluate MetTrain across four primary NLI benchmarks and extensive out-of-distribution (OOD) suites. Overall, MetTrain achieves competitive in-domain performance and strong robustness, surpassing the strongest comparable SSL baseline by up to 2.98 \red{absolute percentage points, abbreviated as pp hereafter,} in MNLI datasets.
	\item \textbf{Enhanced Consistency and Robustness:} Beyond in-domain accuracy, MetTrain improves model reliability under OOD and challenging benchmarks, achieving an improvement of up to 7.00 \red{pp} in OOD Macro-F1 (e.g., when training on SICK and evaluating on SNLI).
\end{itemize}

\section{Background}
\subsection{Natural Language Inference}
Natural Language Inference (NLI)~\cite{DBLP:conf/emnlp/BowmanAPM15} asks a model to identify the logical relationship between a pair of sentences: a ``premise'' and a ``hypothesis''. Table~\ref{tab:nli_examples} illustrates a typical NLI example. These relationships are usually classified as entailment, contradiction, or neutral.

\begin{table}[h]
	\caption{An Example of NLI Task.}
	\label{tab:nli_examples}
	\centering
	\footnotesize
	\renewcommand{\arraystretch}{0.95}
	\setlength{\tabcolsep}{4pt}
	\begin{tabular}{>{\raggedright\arraybackslash}p{0.18\columnwidth}
			>{\raggedright\arraybackslash}p{0.56\columnwidth}
			>{\centering\arraybackslash}p{0.16\columnwidth}}
		\toprule
		\textbf{Item} & \textbf{Content} & \textbf{Label} \\
		\midrule
		Premise & A person in a blue shirt is playing a guitar. & -- \\
		\midrule
		Hypothesis 1 & Someone is playing a musical instrument. & Entailment \\
		Hypothesis 2 & Someone is singing. & Neutral \\
		Hypothesis 3 & No one is playing a guitar. & Contradiction \\
		\bottomrule
	\end{tabular}
\end{table}



NLI has become a standard test for machine reasoning in Natural Language Processing (NLP)~\cite{DBLP:conf/icml/CollobertW08}. 
SOTA models evaluated on SNLI (Stanford NLI corpus) now exceed 90\% accuracy~\cite{DBLP:conf/icann/GajbhiyeMB21}, yet this prowess is inseparable from the sheer volume of human-curated examples that SNLI provides.

Performance collapses on low-resource languages~\cite{conneau-etal-2018-xnli} and specialized test domains~\cite{naik-etal-2018-stress}, and even high-performing models falter on out-of-distribution or adversarial inputs~\cite{hendrycks-etal-2020-pretrained, tian-etal-2021-diagnosing, mccoy-etal-2019-right, gururangan-etal-2018-annotation}. Worse, building new NLI benchmarks remains costly and time-consuming~\cite{nie-etal-2020-adversarial, DBLP:conf/acl/VarshneyBGB22,DBLP:conf/naacl/WilliamsNB18}. To address these challenges of data scarcity and high annotation costs, semi-supervised learning (SSL) emerges as a promising approach~\cite{DBLP:conf/emnlp/ParkC24,DBLP:conf/acl/VarshneyBGB22,DBLP:conf/emnlp/SadatC22a}.


\subsection{Semi-supervised Learning}
Semi-supervised learning (SSL) mitigates data scarcity by leveraging abundant unlabeled text alongside limited labeled examples. However, unlike generic text classification which only requires assigning labels to isolated inputs, NLI additionally demands the construction of Premise-Hypothesis-Label (PHL) triples from unpaired, single-sentence corpora~\cite{DBLP:conf/acl/VarshneyBGB22}, rendering the direct use of unlabeled corpora particularly challenging.

A common strategy for semi-supervised NLI is to synthesize hypotheses from a given premise through rule-based methods~\cite{DBLP:conf/acl/VarshneyBGB22}. Starting from a single premise sentence, such methods apply predefined transformation rules—such as synonym substitution or antonym substitution—to construct candidate hypotheses. While these approaches are easy to implement and computationally efficient, they are inherently limited in their ability to incorporate world knowledge. As a result, they often produce shallow surface-level edits that fail to capture the semantic and inferential diversity required for robust NLI training, leading to marginal and unstable performance gains~\cite{feng-etal-2021-survey,choi-etal-2024-autoaugment,han-etal-2025-rule}. For example, most rule-based methods are constrained to fixed lexical mappings (e.g., synonym swapping), lacking the capacity to model causal consequences, implicit state changes, or higher-order semantic relations.

More recent approaches leverage LLMs to improve semantic richness by exploiting their extensive prior knowledge~\cite{DBLP:conf/emnlp/SadatC22a,DBLP:conf/emnlp/ParkC24}. Given a premise, an often fine-tuned LLM acts as a probabilistic world model to generate one or more hypotheses along with pseudo labels. Although such LLM-based generation enhances fluency and diversity, it comes at a high computational cost~\cite{feng-etal-2021-survey}. Moreover, its effectiveness is coupled to the underlying model's capability~\cite{DBLP:journals/corr/abs-2202-07136}. Relying solely on an unconstrained LLM may introduce hallucinations, semantic drift, or logically inconsistent examples, undermining training reliability~\cite{mckenna-etal-2023-sources}.

%To mitigate these trade-offs, MetTrain treats MRs as explicit, domain-level priors and employs a compact LLM purely as an \emph{MR executor}: MRs specify the inference logic and the resulting labels, while the lightweight LLM is constrained only for realizing these rules in natural language. This decoupling of logical correctness from linguistic realization allows MetTrain to combine the controllability and reliability of rule-based methods with the expressive power of LLMs. As a result, MetTrain generates fluent, diverse, and label-controllable triples, effectively leveraging single-sentence corpora without the prohibitive costs or reliability risks of unconstrained generation.

\subsection{Metamorphic Testing}
Metamorphic Testing (MT) \red{was introduced in the late 1990s and was applied early to numerical analysis programs~\cite{8b7afd955e154a4399cf09e02a74d23e}. It} serves as a validation strategy for scenarios where the correct output is difficult or impossible to determine—a challenge known as the ``oracle problem''~\cite{DBLP:journals/corr/abs-2002-12543}. Instead of verifying individual outputs against a ground truth, MT relies on Metamorphic Relations (MRs) to check the logical consistency between inputs and outputs. \red{MT has been extensively studied and surveyed in the software testing literature~\cite{DBLP:journals/csur/ChenKLPTTZ18,DBLP:journals/tse/SeguraFSC16}, and} has proven effective in validating complex software and AI systems~\cite{DBLP:journals/software/SeguraTZC20,Khokhar2020}.

\red{An MR specifies an expected property between related inputs and their outputs. 
Given a function $f(\cdot)$ and an MR $r = (I_r, O_r)$, a source input $x$ is transformed by $I_r$ into a follow-up input $x' = I_r(x, f(x))$. 
Together, the source and follow-up inputs form a metamorphic group, denoted as $MG_r(x)=\{x,x_r'\}$. 
The corresponding outputs are then expected to satisfy a predefined relation $O_r$, i.e., $O_r(f(x), f(x_r'))$ should hold. 
If this relation is violated, the behavior of $f$ is regarded as faulty.}

\red{For example, in NLI, replacing a word with a synonym can be viewed as a transformation $I_r$. 
If ``happy'' in a hypothesis is replaced with ``joyful'', the semantic meaning should remain unchanged; therefore, the entailment label is expected to be preserved. 
In this case, the original and transformed NLI instances constitute an $MG$, and label preservation is the output relation $O_r$. 
A prediction change from ``entailment'' to ``neutral'' or ``contradiction'' would violate this relation, indicating a potential robustness failure.}

\red{In general MT, the follow-up input may be generated from the source input and, in some cases, its output, i.e., $x' = I_r(x)$ or $x' = I_r(x, f(x))$. In this paper, we focus on the former form of MRs, where the follow-up NLI hypothesis is generated from the source PHL instance and the selected MR, rather than from the target model’s prediction.}

\section{Methodology}

The framework is simple and intuitive. Specifically, MetTrain comprises four steps: (i) pre-checking whether the selected MR or CMR is applicable to the input to avoid structurally incompatible transformations; (ii) using an LLM to generate metamorphic samples under the specified MR or CMR; (iii) performing verification to check format and consistency and filtering out samples that are either excessively or insufficiently transformed, thereby improving sample diversity prior to training; and (iv) training the NLI model on the filtered metamorphic samples together with the original labeled data using the R-Drop loss. Figure~\ref{fig:framework} gives a conceptual overview of MetTrain. 

\red{For a source NLI instance $z$, applying an MR $r$ yields a follow-up instance $z'_r$. 
	The source instance and its follow-up instance together form a metamorphic group, denoted as $MG_r(z)=\{z,z'_r\}$. Since MetTrain uses verified follow-up instances as additional training samples, we refer to the accepted $z'_r$ instances as \emph{metamorphic samples}.}

\begin{figure*}[t]
	\centering
	\includegraphics[width=0.87\linewidth]{Framework}
	\caption{Conceptual Overview of {MetTrain}.}
	\label{fig:framework}
\end{figure*}
%\begin{figure}[t]
%	\centering
%	\includegraphics[width=\linewidth]{SampleGeneration}
%	\caption{An example of metamorphic synthesis from an unlabeled premise.}
%	\label{fig:generation}
%\end{figure}

\subsection{Metamorphic Relation}
MetTrain synthesizes PHL triples using metamorphic relations (MRs) as explicit and verifiable transformation rules. To support both labeled and unlabeled data within a unified formulation, we apply all transformations only to the hypothesis, while keeping the premise unchanged. This design allows the same MR to be used both for augmenting labeled data and for constructing supervision from unlabeled sentences.

Formally, let $\mathcal{T}$ denote natural-language texts, and let
$\mathcal{L}=\{\textsc{Entailment},\textsc{Neutral},\textsc{Contradiction}\}$
denote the label set. An NLI instance is defined as
$z=(P,H,L)\in\mathcal{Z}$, where $\mathcal{Z}=\mathcal{T}\times\mathcal{T}\times\mathcal{L}$,
\red{$P,H\in\mathcal{T}$ is the premise and the hypothesis respectively,} and
$L\in\mathcal{L}$ is the label. \red{We use the same text space $\mathcal{T}$ for premises and hypotheses because both are natural-language sentences.}

\red{\noindent\textbf{Definition 1 (One-source-one-follow-up MR for NLI).}
In this work, we consider a special class of MRs in which each MR maps one
source NLI instance to one follow-up NLI instance. For an MR $r$,
let $\mathcal{D}_r\subseteq\mathcal{Z}$ denote the applicability domain of
$r$, i.e., the set of source instances to which $r$ can be validly applied.
An MR $r$ is specified by a pair of mappings $r=(I_r,O_r)$, where
$I_r:\mathcal{T}\rightarrow\mathcal{T}$ is an input mapping that transforms
the hypothesis according to a predefined rule, and
$O_r:\mathcal{L}\rightarrow\mathcal{L}$ is an output mapping that specifies
the induced label transformation. For every applicable source instance
$z=(P,H,L)\in\mathcal{D}_r$, the corresponding follow-up instance is defined as equation \ref{equ:mr}.
\begin{equation}
	\label{equ:mr}
	z'_r = r(z) = (P, I_r(H), O_r(L))
\end{equation}
Here, $z$ is the source instance and $z'_r$ is the follow-up instance generated
under MR $r$. Thus, the follow-up instance depends only on the source instance
and the selected MR: the premise is kept unchanged, the hypothesis is
transformed by $I_r$, and the label is transformed by $O_r$. This formulation
does not consider MRs that require multiple source instances or simultaneously
generate multiple follow-up instances.}

\red{When $z\notin\mathcal{D}_r$, the MR is regarded as inapplicable and is skipped.
This applicability domain corresponds to the rule-specific pre-checks used in
MetTrain (as shown in Algorithm~\ref{alg:metamorphic_synthesis}). For example, a voice-switch MR is applicable only when the hypothesis
has a suitable syntactic structure, such as subject--verb--object structure.}

For a given unlabeled sentence $S$, we first construct a pseudo-instance
$z=(P,H,L)=(S,S,\textsc{Entailment})$ and then apply applicable MRs to generate
candidate metamorphic samples. In this work, MRs that alter the label are
applied only to entailment instances.

%We define a MR as a pair $r=(I_r,O_r)$, where
%$I_r:\mathcal{T}\rightarrow\mathcal{T}$ transforms the hypothesis according to
%a predefined rule (e.g., synonym or antonym substitution), and
%$O_r:\mathcal{L}\rightarrow\mathcal{L}$ specifies the label mapping induced by
%that transformation. Given an input instance $z=(P,H,L)$, the transformed
%instance under $r$ is defined as
%$z'_r=(P, I_r(H), O_r(L))$.
%For a given unlabeled sentence $S$, we first construct a pseudo-instance $z=(P,H,L)=(S,S,\textsc{Entailment})$ and then apply MRs to generate candidate metamorphic samples.
%In this work, MRs that alter the label are applied only to entailment instances.

%Figure~\ref{fig:generation} presents a simple example of the metamorphic synthesis process for unlabeled data. 

We compile an inventory of 11 elementary MRs in Table~\ref{tab:allMRs}. Certain MRs (e.g., hypernym substitution) are explicitly tailored for unlabeled data to ensure logical validity. 
We group MRs by their label effect:
\(
\mathcal{R}_{\mathrm{inv}}=\{r \mid O_r(L)=L\}, 
\mathcal{R}_{E\rightarrow C}=\{r \mid O_r(E)=C\}, 
\mathcal{R}_{E\rightarrow N}=\{r \mid O_r(E)=N\}.
\)
Here, $L$ stands for the label, while $E,N,C$ represent entailment, neutral, and contradiction, respectively.

\begin{table}[h]
	\centering
	\footnotesize
	\caption{MRs used in MetTrain. Subst. indicates substitution.}
	\label{tab:allMRs}
	\renewcommand{\arraystretch}{0.95}
	\setlength{\tabcolsep}{3pt}
	\begin{tabularx}{\columnwidth}{>{\raggedright\arraybackslash}p{0.24\columnwidth}
			>{\centering\arraybackslash}p{0.12\columnwidth}
			>{\raggedright\arraybackslash}X}
		\toprule
		\textbf{MR Name} & \textbf{Group} & \textbf{Description} \\
		\midrule
		\textsc{Negation Flip} & $\mathcal{R}_{E\rightarrow C}$ & Injects controlled negation into the sentence to induce contradiction. \\
		\textsc{Antonym Subst.} & $\mathcal{R}_{E\rightarrow C}$ & Replaces a keyword with its antonym to reverse the semantic meaning. \\
		\textsc{Add Contradiction} & $\mathcal{R}_{E\rightarrow C}$ & Appends a conflicting phrase or clause to create explicit contradiction. \\
		\textsc{Pronoun Subst.} & $\mathcal{R}_{\mathrm{inv}}$ & Substitutes a noun phrase with an appropriate pronoun while preserving meaning. \\
		\textsc{Voice Switch} & $\mathcal{R}_{\mathrm{inv}}$ & Converts active voice to passive, or vice versa, while maintaining semantic meaning. \\
		\textsc{Synonym Subst.} & $\mathcal{R}_{\mathrm{inv}}$ & Substitutes words with context-appropriate synonyms to diversify vocabulary. \\
		\textsc{Back Translation} & $\mathcal{R}_{\mathrm{inv}}$ & Paraphrases text by translating it into Chinese and back to English. \\
		\textsc{Uninformative} & $\mathcal{R}_{E\rightarrow N}$ & Inserts irrelevant details or context to weaken the entailment relation. \\
		\textsc{Conditional Clause} & $\mathcal{R}_{E\rightarrow N}$ & Extends the hypothesis with conditional or concessive clauses to introduce uncertainty or unrelated information. \\
		\textsc{Extract Snippets} & $\mathcal{R}_{\mathrm{inv}}$ & Extracts a grammatical fragment from the source sentence to form a new entailed hypothesis (\textit{unlabeled data only}). \\
		\textsc{Hypernym Subst.} & $\mathcal{R}_{\mathrm{inv}}$ & Replaces a noun with its hypernym to create a generalized entailment pair (\textit{unlabeled data only}). \\
		\midrule
		\multicolumn{3}{c}{\textit{Composite MRs (CMRs)}} \\
		\midrule
		\textsc{CMR\_Inv} & $\mathcal{R}_{\mathrm{inv}}$ & Randomly combines multiple invariant MRs to create diverse but semantically equivalent hypotheses. \\
		\textsc{CMR\_Flip} & $\mathcal{R}_{E\rightarrow C}$ & Applies label-invariant MRs followed by an $E\!\rightarrow\!C$ MR to produce contradiction. \\
		\textsc{CMR\_Neutral} & $\mathcal{R}_{E\rightarrow N}$ & Mixes invariant MRs with an $E\!\rightarrow\!N$ MR to shift entailment to neutral. \\
		\bottomrule
	\end{tabularx}
\end{table}


We employ the small LLM solely as an MR executor to render the linguistic edits defined by the MRs. We find that a 4B-parameter LLM can generate meaningful and well-formed metamorphic samples. Table~\ref{tab:MRs} presents examples synthesized by Gemma3-4B. Each MR induces a predictable label effect, such as synonym substitution (label-invariant), negation flip ($E\rightarrow C$), or uninformative ($E\rightarrow N$). 




\begin{table*}
	\centering
	\caption{Examples of some MRs used in MetTrain. E, N, C stand for Entailment, Neutral and Contradiction.}
	\label{tab:MRs}
	\resizebox*{0.85\linewidth}{!}{
		\renewcommand{\arraystretch}{0.85} % 缩小行距
		\begin{tabular}{cccc}
			\toprule
			\textbf{MR Type} & \multicolumn{2}{c}{\textbf{Original Pair ($P, H$) and generated Pair ($P, H'$)}} & \textbf{Label}\\
			\midrule
			\multirow{2}{*}{\shortstack{Voice\\Switch}} 
			& $P_1$: The committee approved the proposal. 
			& $H_1$: The committee gave approval to the proposal. 
			& E \\
			& \multicolumn{2}{c}{$H_1'$: The proposal \textcolor{blue}{was approved} by the committee}. 
			& E \\
			\midrule
			\multirow{2}{*}{\shortstack{Uninformative}} 
			& $P_2$: The chef is preparing a complex dish. 
			& $H_2$: The cook is making food. 
			& E \\
			& \multicolumn{2}{c}{$H_2'$: The cook is making food \textcolor{blue}{for a party of 50 guests expected tonight}.} 
			& N \\
			\midrule
			\multirow{2}{*}{\shortstack{Antonym\\Substitution}} 
			& $P_3$: A child is playing in the park. 
			& $H_3$: One young person is enjoying the outdoors. 
			& E \\
			& \multicolumn{2}{c}{$H_3'$: One young person is \textcolor{blue}{avoiding} the outdoors.} 
			& C \\
			\midrule
			\multirow{2}{*}{\shortstack{Hypernym\\Substitution}} 
			& $P_4$: She owns a Siamese cat. 
			& $H_4$: She owns a Siamese cat.
			& E \\
			& \multicolumn{2}{c}{$H_4'$: She owns \textcolor{blue}{an animal}.} 
			& E \\
			\midrule
			\multirow{2}{*}{\shortstack{CMR ($E \rightarrow C$)}} 
			& $P_5$: The report details several key findings. 
			& $H_5$: The file lists multiple important findings. 
			& E \\
			& \multicolumn{2}{c}{$H_5'$: \textcolor{blue}{No important discovery is} reported on the \textcolor{blue}{document, yet the formatting is consistent}.} 
			& C \\
			\bottomrule
		\end{tabular}
	}
\end{table*}

%\paragraph{Composite MRs (CMRs)}
To capture higher-order and more diverse transformations, MetTrain composes MRs into three types of composite MRs (CMRs): (i) \textsc{CMR\_Inv}, which preserves the original label; (ii) \textsc{CMR\_Flip}, which \red{first applies one or more label-invariant MRs to diversify the hypothesis, and then applies an $E\rightarrow C$ MR to produce a contradiction}; and (iii) \textsc{CMR\_Neutral}, which weakens the inference relation without introducing a direct contradiction. The latter two CMRs are applied only to entailment samples.

The composition process is simple and controlled. We first apply \red{one to three randomly selected} label-invariant MRs to increase linguistic diversity, and then append a label-altering MRs to obtain the target label.
Formally, a CMR is defined as an ordered sequence
\red{$c=(r_1,\dots,r_t,\dots,r_m)$}, where $\mathcal{R}_{\mathrm{elem}}$ denotes the set of elementary MRs, $r_t \in \mathcal{R}_{\mathrm{elem}}$ is the MR selected at step $t$, and $m$ is the total number of component MRs in the sequence. Let $z^{(0)} = z = (P, H^{(0)}, L^{(0)})$ denote the original instance, where the superscript $(0)$ indicates the initial state before any transformation is applied. We then apply the component MRs sequentially. For each step $t \in \{1,\dots,m\}$, the intermediate instance is defined as equation \ref{equ:cmr}.
\begin{equation}
	\label{equ:cmr}
	z^{(t)} = (P,H^{(t)},L^{(t)})
	= (P,I_{r_t}(H^{(t-1)}),O_{r_t}(L^{(t-1)}))
\end{equation}
where $z^{(t)}$ is the intermediate instance after step $t$, $H^{(t)}$ and $L^{(t)}$ denote the hypothesis and label after the $t$-th transformation, respectively. The final output of the CMR is $z^{(m)}$.



\subsection{Metamorphic Synthesis}

When using an LLM for data generation, an unconstrained LLM may produce fluent but logically unreliable inference data~\cite{mckenna-etal-2023-sources}, which can degrade the quality of pseudo-labeled training samples~\cite{DBLP:conf/emnlp/ParkC24}. To mitigate this issue, MetTrain uses the LLM not as a free-form pseudo-labeler, but as a constrained MR executor whose transformation effect is predefined by the selected MR. In addition, we apply quality controls in metamorphic synthesis to reduce hallucinations, semantic drift, and malformed outputs during metamorphic synthesis. Algorithm~\ref{alg:metamorphic_synthesis} specifies the dataset-level operational procedure used to construct the metamorphic sample set. Figure~\ref{fig:Generation_Time_Logical_Control_Editable} provides an instance-level illustration of the quality control using a concrete \textsc{Voice Switch} example. 

\begin{figure*}[t]
	\centering
	\includegraphics[width=0.83\linewidth]{Generation_Time_Logical_Control_Editable}
	\caption{Overview of the metamorphic synthesis process, illustrated by \textsc{Voice Switch}.}
	\label{fig:Generation_Time_Logical_Control_Editable}
\end{figure*}

\begin{algorithm}[t]
	\caption{Metamorphic Synthesis}
	\label{alg:metamorphic_synthesis}
	\begin{algorithmic}[1]
		\Require
		\Statex \quad $\mathcal{D}_l \subseteq \mathcal{Z}$ \Comment{labeled PHL triples}
		\Statex \quad  $\mathcal{D}_u \subseteq \mathcal{T}$ \Comment{unlabeled sentences}
		\Statex \quad  $\mathcal{R} = \{r_1,\dots,r_m\}$ \Comment{MR/CMR set}
		\Statex  \quad $A_{\max} \in \mathbb{N}^{+}$ \Comment{maximum retry attempts}
		\Statex  \quad $G$ \Comment{LLM executor}
		\Ensure
		\Statex \quad  $\mathcal{D}_{meta} \subseteq \mathcal{Z}$ \Comment{accepted metamorphic samples}
		
		\State $\mathcal{D}_{meta} \gets \emptyset$ 
		\State $\mathcal{D}_{u}^{\prime} \gets \emptyset$ \Comment{$\mathcal{D}_{u}^{\prime} \subseteq \mathcal{Z}$, pseudo labeled PHL triples}
		
		\ForAll{$S \in \mathcal{D}_u$}
		\State $\mathcal{D}_{u}^{\prime} \gets \mathcal{D}_{u}^{\prime} \cup \{(S,S,\textsc{Entailment})\}$
		\EndFor
		
		\State $\mathcal{D}_{src} \gets \mathcal{D}_l \cup \mathcal{D}_{u}^{\prime}$ \Comment{$\mathcal{D}_{src} \subseteq \mathcal{Z}$, unified source samples}
		
		\ForAll{$z=(P,H,L) \in \mathcal{D}_{src}$} 
		\ForAll{$r \in \mathcal{R}$} 
		\color{red}
		\If{$\Call{PreCheck}{z,r}$}
		
		\cmt{In the previous version, Continue means that the algorithm skips the remaining statements in the current loop body and proceeds to the next iteration. In the previous version, this line indicated that if a sample failed the pre-check, it would be skipped. To avoid confusion, I revised the pseudocode so that generation and verification	are executed only when the sample passes PRECHECK.}
		
		\For{$a \gets 1$ \textbf{to} $A_{\max}$} 
		\State $z_r' \gets \Call{Generate}{G,z,r}$ \Comment{$z_r'\in\mathcal{Z}$}
		\If{$\Call{ConsistencyCheck}{G,z,z_r',r} \land \Call{FormatCheck}{G,z_r',r} $}
		\State $\mathcal{D}_{meta} \gets \mathcal{D}_{meta} \cup \{z_r'\}$
		\State \textbf{break}		
		\EndIf
		\EndFor
		
		\EndIf
		\color{black}
		\EndFor
		\EndFor
		
		\State \Return $\mathcal{D}_{meta}$
		
	\end{algorithmic}
\end{algorithm}

In algorithm~\ref{alg:metamorphic_synthesis}, for unlabeled data, each sentence $S$ is first converted into a pseudo-instance $(S,S,\textsc{Entailment})$, so that both labeled and unlabeled inputs can be processed under the same PHL representation (Line 3--5). The resulting pseudo-labeled set is then merged with the original labeled set to form a unified source pool $\mathcal{D}_{src}$ (Line 6). The subsequent procedure iterates over each source instance $z=(P,H,L)$ and each candidate rule $r\in\mathcal{R}$ to construct metamorphic samples in a controlled manner (Line 7--19). At a high level, Algorithm~\ref{alg:metamorphic_synthesis} defines the corpus-level control flow, whereas Figure~\ref{fig:Generation_Time_Logical_Control_Editable} illustrates the instance-level execution of one synthesis attempt.

$\Call{PreCheck}{ }$ performs rule-specific applicability screening before any generation is attempted (Line 9). Its purpose is to exclude source samples for which the selected MR/CMR would be structurally invalid. Not every input is compatible with every MR/CMR. For example, as illustrated in Figure~\ref{fig:Generation_Time_Logical_Control_Editable}, \textsc{Voice Switch} applies to hypotheses that exhibit a clear subject--verb--object structure and, among other requirements, do not involve intransitive or linking verbs. %Similarly, negation-oriented rules are skipped when the input already contains explicit or implicit negative cues, so as to avoid unintended meaning shifts caused by double negation. 
By filtering out unsuitable cases early, $\Call{PreCheck}{ }$ reduces invalid generations and improves the efficiency of subsequent LLM calls.

$\Call{Generate}{ }$ (Line 11) is guided by constraint-driven prompts that explicitly require semantic faithfulness, preservation of key entities, and adherence to the target MR effect, while restricting the output to a single rewritten hypothesis. This design improves controllability and reduces unnecessary variation in LLM-based augmentation. An example excerpt of the prompt template is shown in Figure~\ref{fig:Generation_Time_Logical_Control_Editable}.

The two verification functions (Line 12--15) in algorithm~\ref{alg:metamorphic_synthesis}, $\Call{ConsistencyCheck}{ }$ and $\Call{FormatCheck}{ }$, are implemented in an MR-specific manner. This verification is designed not only to reject malformed outputs, but also to filter out conversions that are formally valid yet linguistically unnatural or semantically unreliable. An example of these functions on \textsc{Voice Switch} is shown in Figure~\ref{fig:Generation_Time_Logical_Control_Editable}.

$\Call{FormatCheck}{ }$ checks whether the generated hypothesis is usable as a training instance at the surface level. As shown in Figure~\ref{fig:Generation_Time_Logical_Control_Editable}, it first excludes degenerate cases in which the rewritten hypothesis is empty, null, or identical to the original hypothesis, since such outputs do not constitute a valid metamorphic transformation. Second, it filters out a set of manually specified unnatural voice-conversion patterns (e.g., expressions such as ``is had by'' or ``was seemed by''), which are common failure modes when active--passive rewriting is applied to verbs that do not license natural passivization. Third, the candidate hypothesis is further submitted to the LLM for a grammaticality and naturalness judgment, where the model is asked to assess whether the converted sentence is grammatically correct, naturally expressed, and appropriate as a voice transformation. 

$\Call{ConsistencyCheck}{ }$ verifies whether the generated sample preserves the intended MR semantics. For label-invariant rules such as \textsc{Voice Switch}, this means checking whether the transformed hypothesis remains semantically equivalent to the original one and therefore preserves the original entailment relation. The generated hypothesis is compared with the original hypothesis through an LLM-based mutual-entailment check, which tests whether the two sentences bidirectionally entail each other. This step ensures that the transformation is meaning-preserving rather than merely lexically similar. 

A candidate sample is accepted only when both verification functions return true (Line 13). In other words, the output must be simultaneously well formed, grammatically natural, and logically consistent with the transformation effect required by the selected MR. Otherwise, the candidate is rejected and the algorithm retries generation until the retry budget is exhausted.


\subsection{Training Process} 
MetTrain integrates original labeled samples and metamorphic samples within a single training set. R-Drop~\cite{DBLP:conf/nips/LiangWLWMQCZL21} and \red{label smoothing~\cite{7780677}} are introduced to reduce overconfidence and stabilize training under synthetic-label noise. All losses are averaged over the mini-batch. 

For each mini-batch, we perform two forward passes with independent dropout masks, obtaining logits $\mathbf{z}^{(1)}, \mathbf{z}^{(2)} \in \mathbb{R}^{K}$ and predicted class distributions \red{$\mathbf{p}^{(m)}=\mathrm{softmax}(\mathbf{z}^{(m)}), m\in\{1,2\}$, where $\mathrm{softmax}(\mathbf{z})_k=\exp(z_k)/\sum_{j=1}^{K}\exp(z_j)$ for class $k$, and $K$ is the number of classes.} %\cmt{Softmax is a standard operation in neural-network classification, so I initially did not elaborate on it. I have added the formula definition here. I think this is more natural than citing a separate paper for softmax. However, if you think a reference is necessary, I would be happy to add a suitable one.}

To mitigate overfitting, we apply label smoothing with factor $\varepsilon$. For a $K$-class task with $y \in \{1, \dots, K\}$, the smoothed target distribution $\tilde{\mathbf{y}}$ is defined as \red{equation~\ref{equ:smooth_target}. For an original labeled sample, $y$ denotes the ground-truth label; for a metamorphic sample, $y$ denotes the MR-induced label.}
\begin{equation}
	\label{equ:smooth_target}
	\tilde{\mathbf{y}} = (1-\varepsilon) \cdot \mathbf{e}_{y} + \frac{\varepsilon}{K}\mathbf{1},
\end{equation}
where $\mathbf{1} \in \mathbb{R}^{K}$ is an all-ones vector and $\mathbf{e}_{y} \in \mathbb{R}^{K}$ denotes the one-hot vector for class $y$. The $k$-th element is $(\mathbf{e}_{y})_{k} = \mathbb{I}(k=y)$, where $\mathbb{I}(\cdot)$ is the indicator function (equals $1$ if the condition holds, $0$ otherwise).
The supervised cross-entropy (CE) loss is computed as:
\begin{equation}
	\label{equ:CE_def}
	\mathrm{CE}(\mathbf{p}, \tilde{\mathbf{y}}) = -\sum_{k=1}^{K}\tilde{y}_{k}\log p_{k}
\end{equation}
and the CE losses for the two stochastic passes are:
\begin{equation}
	\label{equ:CE}
	\mathcal{L}_{\mathrm{CE}}^{(1)} = \mathrm{CE}(\mathbf{p}^{(1)}, \tilde{\mathbf{y}}), \quad
	\mathcal{L}_{\mathrm{CE}}^{(2)} = \mathrm{CE}(\mathbf{p}^{(2)}, \tilde{\mathbf{y}})
\end{equation}

To \red{enhance} prediction consistency under stochastic dropout, we minimize the symmetric Kullback--Leibler (KL) divergence. The KL divergence between two distributions $\mathbf{p}^{(1)}$ and $\mathbf{p}^{(2)}$ is:
\begin{equation}
	\label{equ:KL_def}
	\mathcal{D}_{\mathrm{KL}}(\mathbf{p}^{(1)}\,\|\,\mathbf{p}^{(2)})=\sum_{k=1}^{K} p^{(1)}_{k}\log\frac{p^{(1)}_{k}}{p^{(2)}_{k}}
\end{equation}
and the symmetric KL term used in R-Drop is defined as:
\begin{equation}
	\label{equ:KL}
	\mathcal{L}_{\mathrm{KL}}^{\mathrm{sym}}
	= \frac{1}{2} \left( \mathcal{D}_{\mathrm{KL}}(\mathbf{p}^{(1)} \,\|\, \mathbf{p}^{(2)})
	+ \mathcal{D}_{\mathrm{KL}}(\mathbf{p}^{(2)} \,\|\, \mathbf{p}^{(1)}) \right)
\end{equation}
\red{The total loss is shown as equation~\ref{equ:loss}. The two CE terms are averaged so that the supervised loss remains on the same scale as a standard single-pass CE loss; the symmetric KL term is then weighted separately by $\lambda$.}
\begin{equation}
	\label{equ:loss}
	\mathcal{L}_{\mathrm{total}}
	= \frac{1}{2}\mathcal{L}_{\mathrm{CE}}^{(1)} + \frac{1}{2}\mathcal{L}_{\mathrm{CE}}^{(2)} + \lambda \mathcal{L}_{\mathrm{KL}}^{\mathrm{sym}}
\end{equation}

\section{Experimental Settings}

\subsection{Research Questions}
We conducted an empirical study to evaluate the following four research questions (RQs):
\begin{itemize}
	\item \red{How is MetTrain's performance compared with those of other SOTA methods?}
	\item How do different MRs affect MetTrain under leave-one-out and cumulative ablation settings?
	\item How does MetTrain perform under distribution shifts and challenging scenarios?
	\item Does MetTrain improve logical consistency and metamorphic validity of NLI models?
\end{itemize}

\subsection{Datasets}

We mainly evaluated MetTrain on four NLI datasets: RTE, SICK, SNLI, and MNLI, \red{with MNLI evaluated on both matched and mismatched development sets}. Our data setup followed standard low-resource practice and is adapted to the scale of each dataset. For RTE and SICK, we used their training splits as labeled data. For SNLI and MNLI, we simulated a low-resource regime by randomly sampling 2,000 labeled examples (class-balanced) and used 2,000 premises of the remaining instances as unlabeled data. Since the test set\red{s} of RTE and MNLI \red{are} not publicly available, we used their \red{respective} development set\red{s} as the test set\red{s} and sampled a small subset
of examples from the training set to be used as the development set.

We generated approximately 10k metamorphic samples for RTE and SICK; around 34k samples for SNLI; and 32k samples for MNLI. To ensure a fair comparison, all competing methods were configured to generate the same number of synthetic samples for each dataset.

%RTE is formulated as a 2-way task (ENTAILMENT vs. NOT\_ENTAILMENT), whereas SICK, SNLI, and MNLI use 3-way labels (ENTAILMENT, CONTRADICTION, NEUTRAL). For OOD evaluation (RQ3) of models trained on RTE, we collapsed CONTRADICTION and NEUTRAL into NOT\_ENTAILMENT when testing on the 3-way datasets. 

In addition to the four main benchmarks, for RQ3 we further assessed robustness on several challenging NLI test sets, including SNLI-hard~\cite{gururangan-etal-2018-annotation}, HANS~\cite{mccoy-etal-2019-right}, DNLI~\cite{welleck-etal-2019-dialogue}, and the NLI stress-test suites~\cite{naik-etal-2018-stress} (antonymy, numerical reasoning, word overlap, negation, length mismatch, and spelling error). These datasets were used only for OOD evaluation.

\subsection{Comparison Methods}

We compared MetTrain against a comprehensive set of baselines, including fine-tuning, prompt-based methods, data augmentation, general SSL, and NLI-specific SSL approaches.

\begin{enumerate}
	
	\item \textbf{Standard Fine-tuning.}
	
	\textbf{BERT:} Fine-tuning a pre-trained BERT model using only labeled data~\cite{DBLP:conf/naacl/DevlinCLT19}, without any unlabeled data or synthetic examples, serving as the fully supervised lower bound. Except for pure prompting baselines (ZSL/ICL), all trainable baselines used bert-base-uncased as the shared backbone.
	
	\item \textbf{Prompt-based Methods.}
	
	\textbf{ICL:} ICL stands for in-context learning on Llama-3-8B-Instruct and Mistral-7B-Instruct. We prepended a fixed class-balanced set of 3 demonstrations per class to each prompt.~\cite{DBLP:conf/nips/BrownMRSKDNSSAA20}
	
	\textbf{ZSL:} ZSL stands for Zero-shot learning on Llama-3.2-3B-Instruct and Llama-3-8B-Instruct without any labeled data~\cite{DBLP:conf/nips/BrownMRSKDNSSAA20}.
	
	\textbf{LM-BFF:} This is a prompt-based fine-tuning method using the manual prompt of Gao et al.~\cite{DBLP:conf/acl/GaoFC20} that reformulates classification as masked language modeling for BERT, using the labeled data only.
	
	\textbf{LM-BFF+Demo:} It extends LM-BFF by further incorporating a few-shot demonstration context into the prompt input.
	
	\item \textbf{Data Augmentation Approaches.}
	
	\textbf{Back Translation:} It generates augmented samples on labeled examples by English-Chinese-English round-trip translation~\cite{edunov2018understanding}.
	
	\textbf{TMix:} It augments labeled training data by randomly selecting pairs of examples and interpolating their hidden representations at a randomly sampled BERT layer from 7, 9 and 12. The corresponding labels are interpolated with the same coefficient~\cite{DBLP:conf/acl/ChenYY20}.
	
	\item \textbf{General SSL methods.}
	
	\textbf{UDA:} It enforces consistency between original and augmented unlabeled samples~\cite{DBLP:conf/nips/XieDHL020}.
	
	\textbf{MixText:} It applies MixUp to labeled and unlabeled data~\cite{DBLP:conf/acl/ChenYY20}.
	
	\textbf{FixMatch:} Pseudo-labeling weakly augmented data and training on strong augmentations~\cite{DBLP:conf/nips/SohnBCZZRCKL20}.
	
	\textbf{FlexMatch:} It extends FixMatch using adaptive confidence thresholds per class~\cite{DBLP:conf/nips/ZhangWHWWOS21}.
	
	\textbf{FreeMatch:} It extends FlexMatch using global/local thresholds with fairness-aware regularization~\cite{DBLP:conf/iclr/Wang0HHFW0SSRS023}.
	
	\textbf{SoftMatch:} It extends FixMatch using a Gaussian function to weight unlabeled samples based on their confidence~\cite{DBLP:conf/iclr/0102TFW0S0RS23}.
	
	\item \textbf{NLI-specific SSL Methods.}
	
	\textbf{PHL-NLI:} Rule-based transformations to generate pseudo-labeled NLI triples~\cite{DBLP:conf/acl/VarshneyBGB22}.
	
	\textbf{SSL4NLI:} It uses LLM-generated pseudo-labeled triples with confidence-based pseudo-label selection and self-debiasing~\cite{DBLP:conf/emnlp/SadatC22a}.
	
	\textbf{VerifyMatch:} It also uses LLM-generated pseudo-labeled triples and further introduces a validator-guided, confidence-aware denoising stage~\cite{DBLP:conf/emnlp/ParkC24}.
	
\end{enumerate}

For In-context learning (ICL) and Zero-shot learning (ZSL) methods, \red{we used the LLM directly as a label predictor. Given a premise--hypothesis pair, the model was asked to determine whether the hypothesis was “True”, “False”, or “Neither” with respect to the premise, corresponding to entailment, contradiction, and neutral, respectively. The prompt template was}: ``\texttt{\{premise\}\textbackslash nQuestion: \{hypothesis\} True, False, or Neither?\textbackslash nAnswer: }'' 

For general SSL methods, \red{the LLM was used to synthesize additional PHL triples from unlabeled premises. Specifically, given an unlabeled premise, the model was prompted to generate one alternative sentence that was definitely true, definitely false, or might be true relative to the original sentence. These three generation targets correspond to entailment, contradiction, and neutral labels, respectively. We used the base Llama-3-8B-Instruct model for general SSL methods, and a fine-tuned Llama-3-8B-Instruct model for SSL4NLI and VerifyMatch.} The specific prompt used for generation was: ``\texttt{We will give you a sentence. Using only the given sentence and your world knowledge, write one alternate sentence that \textbf{is definitely true / is definitely false / might be true} relative to the original. Sentence: premise.}'' 

\cmt{These two paragraphs describe two different uses of LLM prompts in the baseline methods. The first prompt is used for label prediction in ZSL and ICL, where “True”, “False”, and “Neither” correspond to entailment, contradiction, and neutral, respectively. The second prompt is used for data generation in SSL baselines: given an unlabeled premise, the LLM generates an alternative hypothesis that is definitely true, definitely false, or might be true with respect to the premise. These three cases are then mapped to entailment, contradiction, and neutral labels.}

RTE is a binary task. Therefore, for any method that produced three-class outputs, we mapped both CONTRADICTION and NEUTRAL to NOT\_ENTAILMENT on RTE.

For UDA, we adopted back-translation as the augmentation method for both premises and hypotheses. For FixMatch-style SSL methods, we adopted synonym substitution as the weak augmentation and back-translation as the strong augmentation for both premises and hypotheses. 

For any methods that need to generate additional data, we ensured that the number of generated samples matched that of our method for each dataset.

\subsection{Evaluation Metrics}
We used Accuracy and Macro F1 as the primary evaluation metrics. Accuracy is defined as
\begin{equation}
	\text{Acc} = \frac{1}{N}\sum_{i=1}^{N}\mathbb{I}(\hat{y}_i = y_i)
\end{equation}
where $N$ is the total number of test instances, $y_i$ is the ground-truth label of sample $i$, $\hat{y}_i$ is the predicted label, and $\mathbb{I}(\cdot)$ is the indicator function.

Macro F1 computes the mean F1 score across all classes and is adopted here to evaluate how well a model transfers under distribution shifts, especially in cross-dataset OOD settings. For a classification task with $C$ classes, the Macro F1 score is defined as
\begin{equation}
	\text{MacroF1} = \frac{1}{C}\sum_{c=1}^{C} \frac{2\,\text{Prec}_c\,\text{Rec}_c}{\text{Prec}_c+\text{Rec}_c}
\end{equation}
where $C$ is the total number of classes. $\text{Prec}_c$ and $\text{Rec}_c$ are calculated as

\begin{equation}
	\text{Prec}_c = \frac{\text{TP}_c}{\text{TP}_c + \text{FP}_c}, \qquad
	\text{Rec}_c = \frac{\text{TP}_c}{\text{TP}_c + \text{FN}_c}
\end{equation}
where $\text{TP}_c$, $\text{FP}_c$, and $\text{FN}_c$ denote the number of true positives, false positives, and false negatives for each class $c$, respectively.

To evaluate the behavioral consistency of an NLI model under MRs,
we adopted the MR-Violation Rate (MRV)~\cite{9700284}, which quantifies how frequently a model deviates
from the expected relational constraint defined by an MR.  
For a given MR, let $n$ denote the total number of MR-generated test pairs, and let $x$ be the
number of pairs for which the model violates the MR-specified relation (e.g., a
label-invariant MR where the prediction changes, or an entailment-flipping MR where
the label does not flip accordingly). The MRV is defined as
\begin{equation}
	\text{MRV} = \frac{x}{n}
\end{equation}

MRV therefore measures the proportion of test cases in which the model fails to satisfy the
required metamorphic property. Its value lies in $[0,1]$, where lower MRV indicates better
adherence to MR constraints and thus stronger logical consistency and metamorphic validity, while higher
MRV reflects greater deviation from the expected MR behavior.

\subsection{Implementation Details}
We implemented MetTrain in PyTorch (v1.12.1) with HuggingFace Transformers (v4.29.2) on Ubuntu 20.04, using an NVIDIA RTX 3090 GPU (24 GB memory). Except for the pure prompting baselines (ZSL/ICL), all trainable methods employed bert-base-uncased as the NLI backbone, with a maximum sequence length of 128 and a dropout rate of 0.3. For metamorphic synthesis, we used Gemma3-4B as a zero-shot MR executor (without fine-tuning), applying 11 elementary MRs and 3 CMRs to labeled PHL triples and unlabeled premises to generate metamorphic samples.

\red{Hyperparameters were set to commonly used BERT fine-tuning and R-Drop defaults.} The model was trained by minimizing label-smoothed cross-entropy loss ($\varepsilon = 0.1$), regularized by R-Drop with symmetric KL divergence, where $\lambda = 2.0$. Training ran for up to 20 epochs with early stopping, using the AdamW optimizer (learning rate $= 2 \times 10^{-5}$, weight decay $= 10^{-2}$, batch size $= 16$), a linear warmup ratio of 0.06, and gradient clipping with a maximum norm of 1.0. A standard learning-rate scheduler was applied, with the learning rate updated once per optimizer step. The retry budget is set to 5. Final results were averaged over five random seeds; statistical significance was assessed using paired t-tests ($p < 0.05$).

\begin{table*}[h]
	\centering
	\caption{The comparison of accuracy (\%) of MetTrain and other methods. Darker cells indicate better performance in each column. An asterisk indicates a statistically significant difference with the best baseline at $p<0.05$ with paired t-test.}
	\label{tab:comparison}
	\resizebox{0.9\textwidth}{!}{%
		\renewcommand{\arraystretch}{0.8}
		\begin{tabular}{cl|ccccc}
			\toprule
			ID & Methods & \textbf{RTE} & \textbf{SICK} & \textbf{SNLI-2k} & \textbf{MNLI-2k$_m$} & \textbf{MNLI-2k$_{mm}$} \\
			\midrule
			1 & BERT & $66.06 \pm 0.99$ & $85.16 \pm 0.41$ & $74.60 \pm 0.78$ & $63.26 \pm 0.85$ & $64.08 \pm 0.83$ \\
			\midrule
			2 & Llama-3.2-3B ZSL & $48.73 \pm 0.00$ & $51.87 \pm 0.00$ & $40.18 \pm 0.00$ & $35.24 \pm 0.00$ & $34.98 \pm 0.00$ \\
			3 & Llama-3-8B ZSL & $68.95 \pm 0.00$ & $56.60 \pm 0.00$ & $41.40 \pm 0.00$ & $46.82 \pm 0.00$ & $46.33 \pm 0.00$ \\
			4 & Llama-3-8B ICL & $\second{74.36 \pm 0.00}$ & $47.94 \pm 0.00$ & $44.10 \pm 0.00$ & $47.37 \pm 0.00$ & $50.62 \pm 0.00$ \\
			5 & Mistral-7B ICL& \best{$78.33 \pm 0.00$} & $48.08 \pm 0.00$ & $44.30 \pm 0.00$ & $47.88 \pm 0.00$ & $50.55 \pm 0.00$ \\
			6 & LM-BFF & $65.77 \pm 0.78$ & $83.26 \pm 0.88$ & $69.76 \pm 0.82$ & $58.80 \pm 0.90$ & $60.14 \pm 0.87$ \\
			7 & LM-BFF+Demo & $66.71 \pm 1.00$ & $83.98 \pm 1.46$ & $70.23 \pm 0.80$ & $58.49 \pm 0.92$ & $59.22 \pm 0.89$  \\
			\midrule
			8 & Back Trans. & $66.42 \pm 1.38$ & $85.32 \pm 2.40$ & $74.95 \pm 0.75$ & $67.23 \pm 0.78$ & $67.74 \pm 0.76$\\
			9 & TMix & $66.49 \pm 1.54$ & $85.25 \pm 0.91$ & $75.01 \pm 0.74$ & $67.02 \pm 0.79$ & $68.61 \pm 0.73$\\
			\midrule
			10 & UDA & $67.58 \pm 1.78$ & $86.48 \pm 1.38$ & $75.35 \pm 0.72$ & $\third{67.85 \pm 0.71}$ & $68.19 \pm 0.70$ \\
			11 & MixText & $66.14 \pm 1.95$ & $86.50 \pm 1.07$ & $75.84 \pm 0.70$ & $67.57 \pm 0.72$ & $67.77 \pm 0.71$ \\
			\midrule
			12 & FixMatch & $67.87 \pm 1.42$ & $86.47 \pm 1.49$ & $76.11 \pm 0.68$ & $66.77 \pm 0.73$ & $67.53 \pm 0.72$  \\
			13 & FlexMatch & $67.51 \pm 0.84$ & $86.21 \pm 0.68$ & $75.75 \pm 0.69$ & $67.21 \pm 0.71$ & $\third{68.63 \pm 0.69}$  \\
			14 & FreeMatch & $68.23 \pm 0.99$ & $85.98 \pm 0.93$ & $76.03 \pm 0.67$ & $67.48 \pm 0.70$ & $67.92 \pm 0.69$  \\
			15 & SoftMatch & $68.66 \pm 0.78$ & $85.67 \pm 1.22$ & $\third{76.86 \pm 0.65}$ & $67.05 \pm 0.71$ & $68.53 \pm 0.68$\\
			\midrule
			16 & PHL-NLI & $68.88 \pm 1.58$ & $\third{86.58 \pm 1.12}$ & $76.55 \pm 0.65$ & $65.65 \pm 0.65$ & $66.82 \pm 0.74$ \\
			17 & SSL4NLI & $69.10 \pm 2.00$ & $86.22 \pm 0.87$ & $76.38 \pm 0.35$ & $67.91 \pm 0.46$ & $67.63 \pm 0.37$ \\
			18 & VerifyMatch & $70.25 \pm 1.66$ & $\second{86.75 \pm 0.45}$ & $\second{77.96 \pm 0.58}$ & $\second{68.89 \pm 0.65}$ & $\second{69.31 \pm 0.63}$ \\
			\midrule
			19 & \textbf{MetTrain} & $\third{72.71 \pm 1.16}$ & \best{$87.95^* \pm 0.39$} & \best{$78.66^* \pm 0.34$} & \best{$71.18^* \pm 0.35$} & \best{$72.29^* \pm 0.40$} \\
			\bottomrule
		\end{tabular}
	}
\end{table*}


\section{Results and Discussion}
\subsection{RQ1: Performance compared with SOTA}
We evaluated MetTrain against SOTA methods along (1) overall test accuracy, (2) performance under various low-resource settings. Additionally, we discussed the impact of LLM choice on metamorphic synthesis.

\subsubsection{Main Results}
Table~\ref{tab:comparison} reports the performance of MetTrain and the competing methods, presented as mean $\pm$ sample standard deviation. ZSL/ICL are included as reference prompting baselines rather than strictly matched trainable baselines. They were executed only once for efficiency considerations and therefore exhibit zero deviation.

Interestingly, prompt-based LLM baselines are relatively strong on RTE, while remaining much weaker on SICK, SNLI, and MNLI. A plausible reason is that RTE is a small two-class entailment benchmark that prior work has identified as benefiting strongly from transfer learning~\cite{10.5555/3454287.3454581}. In contrast, SICK, SNLI, and MNLI are three-way NLI tasks, which may be harder for prompt-only baselines.

Overall, MetTrain shows consistent, albeit modest, improvements over the comparison methods on most benchmarks. In particular, it yields gains of 2.29 and 2.98 \red{pp} on MNLI-m/mm relative to the strongest SSL baseline. \cmt{The gains on MNLI-m/mm are computed against VerifyMatch: $71.18-68.89=2.29$ pp and $72.29-69.31=2.98$ pp.} These results suggest that metamorphic supervision can serve as a useful auxiliary signal during training by encouraging the model to better respect explicit logical constraints. At the same time, the overall gains remain limited in magnitude, which may be because the proposed MRs primarily encode logical and relational knowledge rather than richer scenario-specific or contextual knowledge. Thus, while MetTrain improves the logical quality of the synthesized supervision, its contribution to in-domain performance is incremental rather than dramatic.

To complement the per-dataset comparison in Table~\ref{tab:comparison}, Figure~\ref{fig:avg_acc} reports the unweighted average accuracy across the five evaluation benchmarks. MetTrain achieves the highest average accuracy of 76.56\%, outperforming the strongest baseline, VerifyMatch \red{(74.63\%)}, by 1.93 \red{pp}. This observation is consistent with the dataset-level results in Table~\ref{tab:comparison}, where MetTrain ranks first on four out of the five benchmarks and remains competitive on RTE. These results suggest that the advantage of MetTrain is not limited to a single dataset, but remains broadly consistent across different evaluation settings. Since the five benchmarks differ in task characteristics and difficulty, this averaged result should be interpreted as a descriptive summary of cross-benchmark consistency rather than a replacement for the detailed per-dataset comparison.

\begin{figure}[t]
	\centering
	\includegraphics[width=0.92\linewidth]{avg_accuracy_across_benchmarks.pdf}
	\caption{Unweighted average accuracy (\%) across the five benchmarks.}
	\label{fig:avg_acc}
\end{figure}


\begin{figure}[!t]
	\centering
	\setlength{\abovecaptionskip}{2pt}
	\setlength{\belowcaptionskip}{-4pt}
	
	\subfloat{%
		\includegraphics[width=0.48\linewidth]{fig_learning_curve_rte}%
		\label{fig:fig1}
	}\hfill
	\subfloat{%
		\includegraphics[width=0.48\linewidth]{fig_learning_curve_sick}%
		\label{fig:fig2}
	}
	
	
	\subfloat{%
		\includegraphics[width=0.48\linewidth]{fig_learning_curve_snli}%
		\label{fig:fig3}
	}\hfill
	\subfloat{%
		\includegraphics[width=0.48\linewidth]{fig_learning_curve_mnli_m}%
		\label{fig:fig4}
	}
	
	
	\subfloat{%
		\includegraphics[width=0.48\linewidth]{fig_learning_curve_mnli_mm}%
		\label{fig:fig5}
	}
	
	\caption{Learning curves showing accuracy under various low-resource settings.}
	\label{fig:learning_curve}
\end{figure}


%\begin{table*}[h]
%	\centering
%	\caption{The comparison of accuracy (\%) of MetTrain with different LLMs in metamorphic synthesis. Dark cells indicate best performance in each column.}
%	\label{tab:LLMablation}
%	\resizebox{0.85\textwidth}{!}{%
%		\renewcommand{\arraystretch}{0.8} % 缩小行距
%		\begin{tabular}{l|ccccc|c}
%			\toprule
%			& \textbf{RTE} & \textbf{SICK} & \textbf{SNLI-2k} &\textbf{MNLI-2k$_m$} &\textbf{MNLI-2k$_{mm}$} &\textbf{Avg.}  \\
%			\midrule
%			MetTrain with Llama-3.2-3B		 		& $72.64 \pm 1.56$ & $87.06 \pm 0.27 $ & $78.41 \pm 0.29$ & {$70.29 \pm 0.59$} & $71.96 \pm 0.46$ & 76.27  \\	
%			MetTrain with Llama-3-8B		 		& $72.06 \pm 1.57$ & $87.37 \pm 0.49$ & $78.39 \pm 0.36$ & $71.08 \pm 0.52$ & {$72.57 \pm 0.44$} & 76.29 \\
%			MetTrain with phi-4 (15B)			 	& $68.38 \pm 2.56$& $87.36 \pm 0.43$ & {$78.25 \pm 0.24$} & \best{$72.06 \pm 0.23$} &  \best{$72.87 \pm 0.18$} & 75.78\\	
%			\midrule
%			MetTrain with Gemma-3-4B		&\best{$72.71 \pm 1.16$} & \best{$87.95 \pm 0.39$} & \best{$78.66 \pm 0.34$} & $71.18 \pm 0.35$ & $72.29 \pm 0.40$ & \best{$76.56$} \\
%			\bottomrule
%		\end{tabular}%
%	}
%\end{table*}



\subsubsection{\red{Varying Low-resource settings}}

\red{We evaluate MetTrain's performance under varying low-resource label budgets.}
Figure~\ref{fig:learning_curve} illustrates how performance changes as the amount of labeled data increases from 500 to 2,000 samples, while the amount of unlabeled data is kept fixed. \red{Although the labeled set increases from 500 to 2,000 examples, all settings remain low-resource compared with the full SNLI/MNLI training sets.}

In the most extreme low-resource setting (500 samples), VerifyMatch shows the strongest initial performance on most datasets, particularly SICK, SNLI, and RTE, likely benefiting from the prior knowledge encoded in its fine-tuned LLM.

Nevertheless, MetTrain remains competitive despite relying on a much smaller, frozen 4B LLM, and generally performs above the standard supervised baseline (BERT). More importantly, as the labeled set grows, MetTrain exhibits a steeper improvement curve than the competing methods. This trend is especially clear on RTE, SNLI, and both MNLI splits, where MetTrain gradually closes the gap and ultimately surpasses VerifyMatch at larger label budgets. On SICK, MetTrain remains highly competitive and stays close to the strongest baseline throughout. Overall, these results suggest that metamorphic supervision is data-efficient, yielding larger benefits as more labeled data becomes available.


\begin{table}[htb]
	\centering
	\caption{The comparison of accuracy (\%) of MetTrain with different LLMs in metamorphic synthesis. Dark cells indicate best performance in each column.}
	\label{tab:LLMablation}
	\resizebox{\columnwidth}{!}{%
		\renewcommand{\arraystretch}{0.8}
		\begin{tabular}{l|ccccc}
			\toprule
			Generator LLM& \textbf{RTE} & \textbf{SICK} & \textbf{SNLI-2k} & \textbf{MNLI-2k$_m$} & \textbf{MNLI-2k$_{mm}$} \\
			\midrule
			Llama-3.2-3B & 72.64 & 87.06 & 78.41 & 70.29 & 71.96  \\
			Llama-3-8B   & 72.06 & 87.37 & 78.39 & 71.08 & 72.57 \\
			Phi-4 (15B)  & 68.38 & 87.36 & 78.25 & \best{72.06} & \best{72.87}  \\
			\midrule
			Gemma-3-4B   & \best{72.71} & \best{87.95} & \best{78.66} & 71.18 & 72.29 \\
			\bottomrule
		\end{tabular}%
	}
\end{table}


\subsubsection{LLM Choice} 
Table~\ref{tab:LLMablation} presents the test accuracy of MetTrain when different LLMs are used for metamorphic synthesis, including Phi, Gemma, and LLaMA variants. Overall, the results are fairly close across models, suggesting that MetTrain is not overly sensitive to the choice of the generator LLM. Among the evaluated models, Gemma-3-4B achieves the best average performance and delivers the strongest results on several benchmarks, although other LLMs remain competitive and in some cases perform best on specific datasets. Notably, no single LLM consistently dominates across all tasks. This pattern indicates that while dataset-specific differences do exist, the effectiveness of MetTrain depends less on scaling up the generator LLM than on the MR-driven supervision itself. The fact that a compact, frozen 4B model can achieve the best overall results further suggests that the gains of MetTrain are primarily attributable to the MR-driven training framework, rather than to the use of a larger generator LLM. 


\begin{tcolorbox}[sharp corners,left=1pt,right=1pt,top=1pt,bottom=1pt,boxsep=1pt,colframe=black]
	\textbf{Answer to RQ1:} MetTrain achieves a competitive performance among 17 comparison methods over 4 datasets.
\end{tcolorbox}



\subsection{RQ2: Ablation analysis of training MRs}
To examine how different MRs affect MetTrain, we conduct two complementary ablation analyses: a leave-one-out study and a cumulative add-one study. Figure~\ref{fig:MRLeaveOneOut} reports the accuracy drop caused by removing a single MR from the full configuration, with the bottom row showing the full-model accuracy. Figure~\ref{fig:MRAblation} presents the cumulative setting, where MR ID 0 denotes the baseline without any MR and IDs 1--14 correspond to progressively adding MRs in a fixed order.
\begin{figure}
	\centering
	\includegraphics[width=1\linewidth]{mr_ablation_heatmap}
	\caption{Leave-one-out MR ablation heatmap for MetTrain, where each row removes one MR from the full configuration. Darker red indicates a larger accuracy (\%) drop after removing the MR. The bottom row presents the full configuration accuracy.}
	\label{fig:MRLeaveOneOut}
\end{figure}

Figure~\ref{fig:MRLeaveOneOut} shows that MR effects are dataset-dependent. On RTE, the largest accuracy drops are observed when removing \textit{pronoun substitution}, \textit{back translation}, and \textit{composite flip} (all $-3.40$ \red{pp}), followed by \textit{hypernym substitution} ($-3.03$). On SICK, the most influential MR is \textit{antonym substitution} ($-1.55$), while \textit{negation flip} and \textit{back translation} also make notable contributions (both $-1.22$). On SNLI, the largest degradations come from removing \textit{back translation} ($-1.42$) and \textit{synonym substitution} ($-1.35$), with \textit{composite invariant} also showing a clear effect ($-1.12$). For MNLI-m/mm, the strongest effect is observed for \textit{composite flip}, which causes drops of $-1.62$ and $-1.51$, respectively; \textit{composite invariant} also yields consistent gains on both splits ($-1.01/-1.00$). Overall, these results do not suggest that any single elementary MR consistently dominates across all datasets. Instead, different benchmarks appear to benefit from different MRs.

Figure~\ref{fig:MRAblation} further illustrates how performance changes as MRs are added cumulatively. Under the fixed addition order used in this study, introducing the first MR (\textit{negation flip}) yields the largest immediate improvement on all five datasets. After this initial increase, the effects of additional elementary MRs become more uneven. RTE shows substantial fluctuations during the elementary-MR stage, whereas SICK and SNLI remain relatively stable after early gains, and both MNLI splits improve more gradually. These patterns suggest that elementary MRs provide complementary supervision, although their effects are not uniformly additive across datasets.

Compared with the elementary-MR stage, the trajectories in the final part of Figure~\ref{fig:MRAblation} become more stable after the introduction of CMRs. This pattern suggests that CMRs may provide additional training signals beyond those offered by individual elementary MRs. However, auxiliary results show that using only CMRs is consistently less effective than using only elementary MRs. \red{Elementary MRs outperform CMR-only training by 1.86--9.54 absolute percentage points, corresponding to approximately 2.67\%--14.11\% relative improvement over the CMR-only setting. Specifically, the relative improvements are 3.85\% on RTE, 4.99\% on SICK, 14.11\% on SNLI, 3.34\% on MNLI-m, and 2.67\% on MNLI-mm.} %65.70 vs.\ 68.23 on RTE, 82.82 vs.\ 86.95 on SICK, 67.62 vs.\ 77.16 on SNLI, 68.02 vs.\ 70.29 on MNLI-m, and 69.72 vs.\ 71.58 on MNLI-mm. 
\red{These results indicate that,} although CMRs are beneficial in the full configuration, they should be better regarded as complementary to elementary MRs rather than as a standalone replacement.

Taken together, the ablation results suggest that elementary MRs and CMRs are complementary rather than substitutive. Elementary MRs account for most of the standalone gains, whereas CMRs appear most useful when combined with elementary MRs, where they provide additional stability and performance improvements.


\begin{tcolorbox}[sharp corners,left=1pt,right=1pt,top=1pt,bottom=1pt,boxsep=1pt,colframe=black]
	\textbf{Answer to RQ2:} The effects of MRs are dataset-dependent. In the ablation settings considered here, elementary MRs provide the main standalone gains, whereas CMRs contribute additional benefits when combined with elementary MRs.
\end{tcolorbox}


\begin{figure}
	\centering
	\includegraphics[width=1\linewidth]{MRAblation}
	\caption{Accuracy (\%) of MetTrain across five NLI datasets under 15 cumulative MR settings. 0 means no MR is applied.}
	\label{fig:MRAblation}
\end{figure}


\subsection{RQ3: Out-of-distribution robustness}


We assess robustness of MetTrain from three complementary perspectives: (1) cross-dataset transfer, (2) challenge-set OOD evaluation, and (3) controlled stress-test robustness.
\subsubsection{Cross-dataset Performance}
We evaluated cross-dataset transfer using Macro F1, as reported in Table~\ref{tab:ood_performance}. When RTE is used, we collapse contradiction and neutral into a single not-entailment class, which may lead to slightly higher absolute scores than those on the original three-way benchmarks. No additional model adaptation or retraining was performed for the OOD evaluation.

\begin{table}[h]
	\centering
	\caption{Out-of-domain Macro F1 comparison. Dark cells indicate best performance.}
	\label{tab:ood_performance}
	\resizebox{1\linewidth}{!}{%
		\renewcommand{\arraystretch}{0.8} % 缩小行距
		\begin{tabular}{c|c|ccccc}
			\toprule
			\textbf{Train Data} & \textbf{Method} & \textbf{RTE} & \textbf{SICK} & \textbf{SNLI} & \textbf{MNLI$_m$} & \textbf{MNLI$_{mm}$} \\
			\midrule
			%            要提一下对于RTE而言，我们把测试数据都改成二分类了，所以指标略高
			\multirow{3}{*}{RTE} & BERT 	& - & 59.06  & 72.38  & 63.57  & 63.85 \\
			& VerifyMatch 					& - & \best{70.83}  & \best{75.22}  & 69.89  & 72.83  \\
			& MetTrain 						& - & {61.74} & {61.03} & \best{74.20} & \best{75.46}\\
			\midrule
			\multirow{3}{*}{SICK} & BERT 	& \best{63.16 } & - & 34.84  & 51.14  & 53.08 \\
			& VerifyMatch 					& 61.98  & - & 40.08  & \best{55.02}  & \best{58.18} \\
			& MetTrain 						& 62.16 & - & \best{47.08} & {50.16} & {52.32}\\
			\midrule
			\multirow{3}{*}{SNLI-2k} & BERT & 55.39  & 40.74  & - & 46.84  & 45.27  \\
			& VerifyMatch 					& 57.08  & 44.73  & - & 57.89  & 58.76  \\
			& MetTrain  					& \best{58.64} & \best{48.52} & - & \best{60.22} & \best{62.25}\\
			\midrule
			\multirow{3}{*}{MNLI-2k} & BERT & 62.04  & 40.93  & 45.48  & - & - \\
			& VerifyMatch 					& 64.58  & 45.27  & 54.95  & - & -  \\
			& MetTrain 						& \best{67.31} & \best{48.26} & \best{57.84} & - & - \\
			\bottomrule
		\end{tabular}%	
	}
\end{table}

Overall, the results suggest that MetTrain provides stronger cross-dataset transfer in several, though not all, train$\rightarrow$test settings. Its advantage is most pronounced when the source data are SNLI-2k or MNLI-2k. In these cases, MetTrain consistently achieves the best Macro F1 on all target datasets. When trained on SICK, MetTrain markedly improves transfer to SNLI by up to 7.00 \red{pp}, but remains behind BERT on RTE and behind VerifyMatch on MNLI$_m$ and MNLI$_{mm}$. Likewise, under RTE$\rightarrow$target transfer, MetTrain performs best on the two MNLI splits, whereas VerifyMatch is stronger on SICK and SNLI. Taken together, these results indicate that metamorphic supervision can improve OOD generalization, especially when learned from larger and more diverse source benchmarks, although the benefit is not uniform across all source-target combinations.

\subsubsection{Challenging OOD Performance}

To further evaluate the robustness of MetTrain, we considered more challenging OOD benchmarks in Table~\ref{tab:ood_hard}, where models trained on SNLI-2k or MNLI-2k were tested on SNLI-hard~\cite{gururangan-etal-2018-annotation}, HANS~\cite{mccoy-etal-2019-right}, and DNLI~\cite{welleck-etal-2019-dialogue}. Overall, MetTrain shows \red{slight but consistent improvements in most of} these harder transfer settings. \red{Because some absolute differences are small, we treat these results as supportive rather than conclusive.} When trained on SNLI-2k, it achieves the best accuracy on SNLI-hard and DNLI, while remaining highly competitive on HANS, where its result is only marginally below VerifyMatch. When trained on MNLI-2k, MetTrain obtains the \red{numerically} best performance on all three benchmarks. Its advantage is particularly clear on DNLI, where it outperformed VerifyMatch from 35.27 to 37.09 under SNLI-2k training and from 45.28 to 46.90 under MNLI-2k training. These results \red{provide supportive evidence that} metamorphic supervision \red{can} improves robustness on more challenging OOD benchmarks and may help the model generalize better to adversarial and dialogue-oriented NLI settings.


\begin{table}[h]
	\centering
	\caption{The comparison of accuracy (\%) of MetTrain and other methods on challenging out-of-distribution	data. Dark cells indicate best performance. \red{An asterisk indicates a statistically significant difference with the best baseline at $p<0.05$ with paired t-test.}}
	\label{tab:ood_hard}
	\resizebox{\linewidth}{!}{%
		\renewcommand{\arraystretch}{1} % 缩小行距
		\begin{tabular}{l|ccc|ccc}
			\toprule
			Train data & \multicolumn{3}{c|}{\textbf{SNLI-2k}} & \multicolumn{3}{c}{\textbf{MNLI-2k}}\\
			Test data & \textbf{SNLI-hard} & \textbf{HANS} & \textbf{DNLI} &\textbf{SNLI-hard} & \textbf{HANS} & \textbf{DNLI}  \\
			\midrule
			BERT 				& 61.12 & 49.59 & 32.39 & 42.16 & 49.81 & 36.70 \\
			VerifyMatch			& {63.30} & \best{50.06} & 35.27 & 49.27 & 50.31 & 45.28 \\
			MetTrain		 	& \best{63.84} & {50.00} & \best{37.09\red{$^*$}} & \best{49.94} & \best{50.34} & \best{46.90\red{$^*$}} \\
			\bottomrule
		\end{tabular}%
	}
\end{table}


\begin{figure}[!t]
	\centering
	\subfloat{%
		\includegraphics[width=0.54\linewidth]{radar_chart_RTE}%
	}
	\subfloat{%
		\includegraphics[width=0.49\linewidth]{radar_chart_SICK}%
	}
	
	\subfloat{%
		\includegraphics[width=0.49\linewidth]{radar_chart_SNLI}%
	}
	\subfloat{%
		\includegraphics[width=0.49\linewidth]{radar_chart_MNLI}%
	}
	\caption{Radar charts comparing the robustness of MetTrain and other methods on challenging OOD evaluation benchmarks (NLI stress test). The axes represent accuracy (\%) across 11 distinct linguistic challenges.}
	\label{fig:stress_test}
\end{figure}


\subsubsection{NLI Stress Test Performance}
Figure~\ref{fig:stress_test} further probes robustness on a suite of NLI stress tests~\cite{naik-etal-2018-stress}, covering competence (antonymy, numerical reasoning), distraction (word overlap, negation, length mismatch), and noise (spelling errors). 

The results reveal a selective rather than uniform robustness pattern. Across all four training corpora, MetTrain consistently performs better on antonymy and spelling-error conditions, and it also shows clear advantages on most length-mismatch settings. For example, when trained on MNLI, MetTrain surpasses  VerifyMatch in antonymy from 27.97\%/29.98\% (m/mm) to 46.65\%/45.52\%, and further raises length-mismatch accuracy to 66.02\%/68.25\% and spelling-error accuracy to 63.49\%/64.21\%. Similar trends appear on the other datasets. MetTrain achieves the highest antonymy scores and spelling-error results across all four training corpora.

One possible explanation for these gains is the metamorphic supervision. MRs such as \textit{antonym substitution} and \textit{negation flip} create minimally different premise--hypothesis pairs in which a single antonym or negation must invert the label. This encourages the model to treat such cues as decisive, thereby reducing errors when antonymic changes occur. Neutralizing MRs---such as \textit{uninformative} and \textit{conditional clause}---inject irrelevant or hedging information into hypotheses, often making them longer. This may indirectly reduce the model's sensitivity to sentence-length variation, which helps explain the gains on length-mismatch conditions. At the same time, these MRs encourage the model to focus more on core semantic content and less on semantically uninformative surface variations, which may also contribute to the improvement under spelling-error conditions, even though they do not explicitly introduce spelling noise.

However, MetTrain does not excel on all stress dimensions. Its performance is generally weaker on numerical reasoning and word overlap, where VerifyMatch or even BERT remains stronger on several datasets. This is especially visible on RTE, SICK, and MNLI for word overlap and numerical reasoning. This pattern mirrors our current MR inventory (Table~\ref{tab:allMRs}), which primarily targets coarse semantic and syntactic transformations but lacks explicit operations over numbers, quantities, or overlap-based lexical heuristics, leaving arithmetic, comparison, and overlap-sensitive reasoning insufficiently covered.

Overall, the stress-test results suggest that MetTrain mainly improves robustness to antonymic perturbations, spelling noise, and sentence-length distractions, while leaving more room for improvement on numerical reasoning and word-overlap-heavy cases. These findings indicate that MR-driven supervision can mitigate several common shallow heuristics, but its benefits remain dependent on the types of transformations encoded in the current MR set.

\begin{tcolorbox}[sharp corners,left=1pt,right=1pt,top=1pt,bottom=1pt,boxsep=1pt,colframe=black]
	\textbf{Answer to RQ3:} MetTrain shows improved robustness under cross-dataset transfer and challenging OOD benchmarks, and demonstrates gains on certain stress-test conditions.
\end{tcolorbox}


\subsection{RQ4: Logical Consistency and Metamorphic Validity}
To answer RQ4, we investigate the logical consistency and metamorphic validity of MetTrain through metamorphic testing and qualitative case analysis.
\subsubsection{Metamorphic testing}
We conduct metamorphic testing based on a set of 11 elementary MRs. Unlike prior works that rely on a fixed split of MRs, we adopt a randomized split validation strategy to ensure a more rigorous and unbiased evaluation. Specifically, we partition the inventory of 11 elementary MRs into two disjoint subsets: a training set (6 MRs) used to generate supervision signals, and a held-out testing set (5 MRs) which is applied exclusively to the test data for evaluation. We report the final MR Violation Rates (MRV) on the held-out testing MRs. This design allows us to assess the generalization of logical consistency to unseen MRs.

\paragraph{Results} 
Table~\ref{tab:rq4_mrv_unseen} reports the MRV on test MRs across different datasets under randomized splits. Overall, MetTrain consistently achieves lower MRV, indicating improved logical consistency and generalization. In particular, substantial reductions are observed on SICK and SNLI, where the violation rates decrease from $54.43\%$ to $41.69\%$ and from $52.52\%$ to $37.89\%$, respectively. MetTrain also outperforms BERT on MNLI$_{mm}$, reducing the MRV from $55.68\%$ to $48.24\%$, and achieves a slight improvement on MNLI$_m$. Although the performance gain on RTE is relatively modest, MetTrain still yields a statistically significant reduction in MRV. Overall, these results demonstrate the effectiveness of MetTrain in mitigating MR violations and enhancing logical robustness across diverse NLI benchmarks.


\begin{table*}[t]
	\centering
	\caption{MR Violation Rates (MRV, \%) on test MRs under randomized splits. Lower MRV indicates better logical consistency and generalization. An asterisk indicates a statistically significant difference with the baseline at $p<0.05$ with paired t-test.}
	\label{tab:rq4_mrv_unseen}
	\resizebox{\linewidth}{!}{%
		\renewcommand{\arraystretch}{0.8} 
		\begin{tabular}{lccccc}
			\toprule
			& \multicolumn{5}{c}{\textbf{MR Violation Rates (MRV) $\downarrow$}} \\
			\cmidrule(lr){2-6}
			& \textbf{RTE} & \textbf{SICK} & \textbf{SNLI} & \textbf{MNLI$_m$} & \textbf{MNLI$_{mm}$} \\
			\midrule
			BERT & $51.21 \pm 1.23$ & $54.43 \pm 6.92$ & $52.52 \pm 1.65$ & $52.70 \pm 2.17$& $55.68 \pm 1.35$ \\
			{MetTrain} & ${49.56^* \pm 1.08}$ & ${41.69^* \pm 2.66}$ & ${37.89^* \pm 10.22}$ & ${51.62 \pm 10.64}$ & ${48.24^* \pm 5.87}$\\
			\midrule
			MRV Reduction (BERT $-$ MetTrain) ($\uparrow$) & $1.65$ & $12.74$ & $14.63$ & $1.08$ & $7.44$ \\
			\bottomrule
		\end{tabular}%
	}
\end{table*}

\begin{table*}[t]
	\centering
	\caption{Illustrative MR violation cases of BERT. Predictions that \emph{violate} the MR are underlined.}
	\label{tab:mrv-cases}
	\resizebox{\linewidth}{!}{%
		\renewcommand{\arraystretch}{1.0}
		\begin{tabular}{p{1.4cm} p{4.3cm} p{5cm} p{1.15cm} p{0.6cm} p{1.2cm}}
			\toprule
			MR & Premise & Hypothesis (orig. $\rightarrow$ MR-transformed) &
			Label relation & BERT  & MetTrain \\
			Antonymy substitution &
			\textit{In the other bracket, the Broncos beat the New York Jets.} &
			orig: \textit{The Broncos beat the New York Jets.} \newline
			$\rightarrow$ \textit{Nobody beats the New York Jets.} &
			E$\rightarrow$\textbf{C} &
			\underline{\textbf{E}} &
			\textbf{C} \\
			\midrule
			Voice switch &
			\textit{From the corner of his eye he saw Jamus look over the broken mare.} &
			orig: \textit{Jamus looked over the mare.} \newline
			$\rightarrow$ \textit{The mare was looked over by Jamus.} &
			E$\rightarrow$\textbf{E} &
			\underline{\textbf{C}} &
			\textbf{E} \\
			\midrule
			Uninfor- mative &
			\textit{US President Barack Obama has been sketching out the future direction of his administration, in events to mark his first 100 days in office $\cdots$ (73 words)} &
			orig: \textit{Barack Obama has been President of the U.S. for 100 days.} \newline
			$\rightarrow$ \textit{Barack Obama has been President of the U.S. for 100 days while wearing his favorite tie.} &
			E$\rightarrow$\textbf{N} &
			\underline{\textbf{E}} &
			\textbf{N} \\
			\bottomrule
		\end{tabular}
	}
\end{table*}

\subsubsection{Case studies of MR violations}
While MRV provides an aggregate measure of logical consistency, we further conduct a
qualitative analysis of representative MR violations before and after MetTrain.
Table~\ref{tab:mrv-cases} presents three illustrative cases covering different types of MRs, \textit{antonymy substitution}, \textit{voice switch}, and \textit{uninformative}.

These examples in Table~\ref{tab:mrv-cases} reflect systematic weaknesses of the BERT baseline rather than isolated errors.
In our MR test logs, BERT generally shows high MRV on some MR families: for instance,
\textit{antonymy substitution} reaches up to $87.30\%$ violation rate (873/1000) in SNLI, and
\textit{voice switch} yields up to $67.74\%$ violation rate (147/217) in MNLI$_m$, suggesting that the model often struggles with both decisive lexical cues and meaning-preserving syntactic changes.

Across the three cases shown in Table~\ref{tab:mrv-cases}, we observe two recurring failure modes: (i) over-sensitivity
to format variations that should preserve the label, and (ii) under-sensitivity to small
semantic edits that should change the label. 

As shown in Table~\ref{tab:mrv-cases}, for \textit{antonymy substitution},
a single cue (e.g., ``Nobody beats $\cdots$'') should deterministically flip entailment into
contradiction. BERT fails to flip, but MetTrain responds correctly. For \textit{voice switch}, the transformed hypothesis is a passive paraphrase of the original entailed statement, yet BERT incorrectly predicts contradiction, whereas MetTrain maintains the expected entailment. Finally, for the MR \textit{uninformative}, adding unsupported details should weaken entailment to neutral. However, BERT remains over-confident in entailment, while MetTrain produces the MR-consistent neutral label.




%Both BERT and VerifyMatch violated even simple
%invariant MRs: introducing a synonym (\emph{playing a song}) in the hypothesis can spuriously
%change predictions from entailment to neutral. Under label-flipping MRs such as \textsc{negation\_flip},
%they often fail to invert the label when a targeted negation is added. In contrast, MetTrain respects
%both invariance and flip constraints in these cases. Moreover, despite never seeing quantifier-based
%MRs during training, MetTrain also better handles the unseen monotonicity transformation
%(\emph{every} $\rightarrow$ \emph{some}), suggesting that MR-guided supervision induces transferable
%logical regularities beyond the exact MRs used for training.

\begin{tcolorbox}[sharp corners,left=1pt,right=1pt,top=1pt,bottom=1pt,boxsep=1pt,colframe=black]
	\textbf{Answer to RQ4:} Reduced MRV under randomized MR splits indicates that MetTrain improves logical consistency.
\end{tcolorbox}

\section{Threats to Validity}

\subsection{Construct validity.}
MetTrain relies on a manually curated set of 11 elementary MRs and 3 composite MRs. Although these MRs are grounded in prior work on metamorphic testing and NLI, they cover only a limited portion of the possible MR space. The 14 MRs used in this study should therefore be viewed as a constrained subset rather than a complete characterization of the MR design space. Recent work has catalogued nearly 200 MRs for LLM-based NLP tasks~\cite{11185922}; incorporating additional relations may further benefit training.

Another threat concerns the interpretation of MR ablation results. The effect of a given MR may vary across datasets, and its contribution may depend on which other MRs are present. Therefore, the reported ablations provide only an empirical view of MR effects under the studied settings, rather than a complete account of all MR contributions and interactions. More systematic attribution strategies and richer MR combinations remain future work.

\subsection{Internal validity.}
We use a single 4B LLM to instantiate and verify MRs, so generation errors or imperfect transformations may influence the final results. Although pre-checks and verification help reduce such risks, they cannot eliminate them completely. Using a separate verifier or stronger cross-model validation may further improve reliability.


\subsection{External validity.}
Our evaluation covers four English NLI benchmarks together with standard OOD and stress-test datasets. Although these settings span a range of inference scenarios, the results may not directly generalize to other languages, specialized domains (e.g., biomedical or legal NLI), or other reasoning tasks. Extending MetTrain to such settings may require not only a larger MR inventory, but also more task-specific and language-aware MR designs beyond the mainly lexical and syntactic relations considered in this work.

%\subsection{Conclusion validity.}
%Although all results are averaged over five random seeds and statistical significance is assessed with paired t-tests, the number of runs remains limited, especially for small datasets such as RTE, where performance may vary more across splits and random initialization. In addition, some conclusions drawn from the ablation study should be interpreted with caution, since MR effects may depend on the fixed addition order in cumulative analysis and on interactions among MRs in leave-one-out analysis. Therefore, while the results consistently support the effectiveness of MetTrain, the magnitude of individual MR contributions should not be over-interpreted.




\section{Related Work}
\subsection{Metamorphic Testing in Software and AI}
Metamorphic testing (MT) has a rich history in software engineering as a strategy to ``test the untestable''~\cite{DBLP:journals/software/SeguraTZC20}. By defining transformations that any correct program should uphold, MT enables the generation of follow-up test cases without a traditional oracle. It has now been embraced across various AI fields. For example, Xie et al.~\cite{XIE2011544} tested machine-learning classifiers via MRs, showing that even well-performing models can exhibit numerous MR violations, and follow-up work has used MT to assess AI systems, including machine translation~\cite{DBLP:conf/icse/WangHWZHLHL23, DBLP:journals/tosem/XieJCC24}, image captioning~\cite{DBLP:journals/tse/XieLC24, DBLP:conf/kbse/YuanPW22}, and LLM fairness~\cite{DBLP:conf/qrs/Li0CXG24, 9700284, ribeiro-etal-2020-beyond}. The common theme is \emph{testing}---detect violations \emph{after} training. This trend highlights a missed opportunity: if MRs can detect flaws, could they prevent them in the first place?


A smaller but growing thread explores this by injecting MRs \emph{during} training to enhance robustness or accuracy~\cite{10.1145/3679006.3685067,DBLP:journals/corr/abs-2412-01958,DBLP:conf/icse/XuTFBZC18,DBLP:journals/corr/abs-2104-04718}. 
While promising, this shift from evaluation to supervision is still in its infancy. Initial studies have primarily focused on domains like image classification~\cite{10.1145/3679006.3685067,DBLP:journals/corr/abs-2412-01958,DBLP:conf/icse/XuTFBZC18,DBLP:journals/corr/abs-2104-04718}, where MRs often involve intuitive geometric or color-space transformations. Whether these principles can be extended to tasks demanding high-level logical reasoning and semantic fidelity, such as NLI, remains a critical open question. This calls for a systematic investigation into the feasibility and effectiveness of using MRs for training in complex, logic-intensive domains.

\subsection{Semi-supervised NLI}
Semi-supervised learning (SSL) for NLI seeks to reduce dependence on expensive labeled corpora such as SNLI by exploiting large amounts of unlabeled text. Generic SSL methods—including consistency-based approaches like Unsupervised Data Augmentation (UDA)~\cite{DBLP:conf/nips/XieDHL020} and interpolation based methods like MixText~\cite{DBLP:conf/acl/ChenYY20} train models to produce stable predictions under data augmentation, often using back-translation~\cite{edunov2018understanding} or hidden-space mixup~\cite{zhang2017mixup}.  Pseudo-labeling-driven algorithms such as FixMatch~\cite{DBLP:conf/nips/SohnBCZZRCKL20} and SoftMatch~\cite{DBLP:conf/iclr/0102TFW0S0RS23} further combine confidence-filtered pseudo labels with strong augmentation and sample-weighting strategies to better exploit unlabeled data. Although these techniques can be applied to NLI by treating each premise–hypothesis pair as a generic text input, they largely ignore the logical structure of NLI labels. 

NLI-specific SSL work goes a step further: Sadat et al.~\cite{DBLP:conf/emnlp/SadatC22a} propose a semi-supervised framework (SSL4NLI) that uses confidence-aware pseudo-labels and self-debiased training to improve robustness, especially under annotation artifacts.  Varshney et al.~\cite{DBLP:conf/acl/VarshneyBGB22} (PHL-NLI) generate premise-hypothesis-label triples from unlabeled text in a fully unsupervised way and train NLI models solely on such synthetic data.  Most recently, VerifyMatch~\cite{DBLP:conf/emnlp/ParkC24} couples an 8B-parameter LLM teacher with a learned validator and confidence-aware MixUp to filter noisy pseudo-labels, achieving strong in-domain and cross-domain performance but at non-trivial computational cost and without explicit guarantees of logical consistency.  

Overall, these methods demonstrate the benefit of unlabeled data for NLI, yet they treat logical consistency as an emergent property of the model, rather than encoding explicit, checkable relational constraints into the training objective.


\section{Conclusion}

In this paper, we proposed \textbf{MetTrain}, a framework that shifts the role of metamorphic relations from passive post-hoc test oracles to the source of training supervision. By decoupling logical correctness from linguistic generation, MetTrain leverages a lightweight, frozen 4B-parameter LLM to enforce 11 types of logical constraints, effectively mitigating the noise and semantic drift inherent in traditional semi-supervised learning.

This shift from data-driven supervision to property-driven supervision shows great potential. Evaluations across four benchmarks confirm that this MR-driven framework delivers competitive and often better performance compared to state-of-the-art approaches. Furthermore, it demonstrates enhanced robustness on out-of-distribution datasets, stress tests, and metamorphic testing. Our work demonstrates that embedding explicit, testable software properties into the training loop is essential for constructing trustworthy and logically consistent AI systems.


\section*{Acknowledgments}
This work was partly supported by the National Natural Science
Foundation of China (NSFC) (Grant nos. 62172194, 62202206 and
U1836116), Qinglan Project of Jiangsu Province, and the Graduate Research
Innovation Project of Jiangsu Province (Grant number:
KYCX26\_4238).



%{\appendix[Proof of the Zonklar Equations]
%Use $\backslash${\tt{appendix}} if you have a single appendix:
%Do not use $\backslash${\tt{section}} anymore after $\backslash${\tt{appendix}}, only $\backslash${\tt{section*}}.
%If you have multiple appendixes use $\backslash${\tt{appendices}} then use $\backslash${\tt{section}} to start each appendix.
%You must declare a $\backslash${\tt{section}} before using any $\backslash${\tt{subsection}} or using $\backslash${\tt{label}} ($\backslash${\tt{appendices}} by itself
% starts a section numbered zero.)}



%{\appendices
%\section*{Proof of the First Zonklar Equation}
%Appendix one text goes here.
% You can choose not to have a title for an appendix if you want by leaving the argument blank
%\section*{Proof of the Second Zonklar Equation}
%Appendix two text goes here.}



%\section{References Section}
%You can use a bibliography generated by BibTeX as a .bbl file.
% BibTeX documentation can be easily obtained at:
% http://mirror.ctan.org/biblio/bibtex/contrib/doc/
% The IEEEtran BibTeX style support page is:
% http://www.michaelshell.org/tex/ieeetran/bibtex/
 
 % argument is your BibTeX string definitions and bibliography database(s)
%\bibliography{IEEEabrv,../bib/paper}
%
%\section{Simple References}
%You can manually copy in the resultant .bbl file and set second argument of $\backslash${\tt{begin}} to the number of references
% (used to reserve space for the reference number labels box).
% 
\bibliographystyle{IEEEtran}
\bibliography{sample-base.bib}
%
%\newpage
%
%\section{Biography Section}
%If you have an EPS/PDF photo (graphicx package needed), extra braces are
% needed around the contents of the optional argument to biography to prevent
% the LaTeX parser from getting confused when it sees the complicated
% $\backslash${\tt{includegraphics}} command within an optional argument. (You can create
% your own custom macro containing the $\backslash${\tt{includegraphics}} command to make things
% simpler here.)
% 
%\vspace{11pt}
%
%\bf{If you include a photo:}\vspace{-33pt}
%\begin{IEEEbiography}[{\includegraphics[width=1in,height=1.25in,clip,keepaspectratio]{fig1}}]{Michael Shell}
%Use $\backslash${\tt{begin\{IEEEbiography\}}} and then for the 1st argument use $\backslash${\tt{includegraphics}} to declare and link the author photo.
%Use the author name as the 3rd argument followed by the biography text.
%\end{IEEEbiography}
%
%\vspace{11pt}
%
%\bf{If you will not include a photo:}\vspace{-33pt}
%\begin{IEEEbiographynophoto}{John Doe}
%Use $\backslash${\tt{begin\{IEEEbiographynophoto\}}} and the author name as the argument followed by the biography text.
%\end{IEEEbiographynophoto}
%
%
%
%
%\vfill

\end{document}


